%!TeX program = xelatex
\documentclass{SYSUReport}
\usepackage{amsmath}
\usepackage{tabularx}
\usepackage{float}
\usepackage{multirow}
\usepackage{graphicx}

\headl{}
\headc{}
\headr{并行程序设计与算法实验}
\lessonTitle{并行程序设计与算法实验}
\reportTitle{Lab7-MPI并行应用}
\stuname{林宏宇}
\stuid{22320001}
\inst{计算机学院}
\major{计算机科学与技术}
\date{2026年5月18日}

\begin{document}

\cover
\thispagestyle{empty}
\clearpage

\section{实验目的}
\begin{itemize}
   \item 验证并行FFT的加速效果与可扩展性
   \item 评估数据打包对通信性能的优化作用
   \item 分析并行规模对内存消耗的影响
\end{itemize}

\section{实验内容}
\begin{itemize}
    \item \textbf{串行FFT分析与并行化改造：}
    \begin{itemize}
        \item 阅读并理解提供的串行傅里叶变换代码 (fft\_serial.cpp)。
        \item 使用MPI对串行FFT代码进行并行化改造。
    \end{itemize}
    \item \textbf{MPI数据通信优化：}
    \begin{itemize}
        \item 应用MPI数据打包技术（\texttt{MPI\_Pack}/\texttt{MPI\_Unpack}）来对通信数据进行重组，以优化消息传递效率。
    \end{itemize}
    \item \textbf{程序性能与内存分析：}
    \begin{itemize}
        \item 分析在不同并行规模及不同问题规模条件下，并行FFT程序的加速比和并行效率。
        \item 分析数据打包技术对于并行程序整体性能的具体影响。
        \item 使用Valgrind Massif工具采集并分析并行程序在不同配置下的内存消耗。
    \end{itemize}
\end{itemize}

\section{实验环境}
\begin{itemize}
    \item \textbf{CPU：}Intel Xeon Gold 6240R @ 2.40 GHz（最大睿频 4.00 GHz），Cascade Lake 微架构，支持 AVX-512 指令集
    \item \textbf{核心与线程：}2 路 NUMA 节点（NUMA node0: CPU 0--23, 48--71；NUMA node1: CPU 24--47, 72--95），每路 24 核 48 线程，共 48 物理核 / 96 逻辑线程（超线程开启），本次实验仅使用物理核范围内的进程数（$P \leq 16$），未触及超线程
    \item \textbf{Cache 层次：}L1d/L1i: 32 KB/核（共 1.5 MB）；L2: 1024 KB/核（共 48 MB）；L3: 35.75 MB/NUMA 节点（共 71.5 MB）
\end{itemize}

\section{实验结果与分析}

\subsection{并行FFT的实现与正确性验证}

\textbf{算法设计：}

本实验基于Cooley-Tukey按频率抽取（DIF）蝶形算法实现并行FFT。串行核心函数为\texttt{step}，该函数在每一阶段（参数$mj = 1, 2, 4, \dots, N/2$）执行蝶形运算：将输入数组分为前半部分$a$（全局复位置 $0 \dots N/2-1$）和后半部分$b$（全局复位置 $N/2 \dots N-1$），在每个蝶形组内计算$a+b$（和值）与$W \cdot (a-b)$（差值乘以旋转因子），分别输出至结果数组的上下两部分。

\textbf{并行化策略：}

采用一维块状数据划分。$N$个复数等分为$P$份，进程$r$持有$N/P$个连续的复数（即$2N/P$个双精度浮点数），覆盖全局复位置$[r \cdot N/P,\; (r+1) \cdot N/P)$。对于所有FFT阶段，前半部分数据始终分布在进程$0 \sim P/2-1$，后半部分数据始终分布在进程$P/2 \sim P-1$。

每阶段通信与计算的三步流程如下：
\begin{enumerate}
    \item \textbf{数据交换：}进程$r$与伙伴进程$r \oplus (P/2)$通过MPI\_Sendrecv交换其持有的全部$N/P$个复数。交换后双方均拥有对应$a$和$b$位置上的完整蝶形输入对。
    \item \textbf{局部计算：}每个进程遍历其$a$-位置集合（共$N/P$个，P=1特殊处理为$N/2$个），对每个位置计算蝶形运算，得到和值与差值两个输出值。
    \item \textbf{输出重分布：}使用MPI\_Alltoallv将计算得到的输出值按目标全局位置发送至正确进程，确保输出数组保持块状分布以供下一阶段使用。
\end{enumerate}

该方案要求$P$为2的幂且$N/P$为整数。P=1时无需进程间通信，直接使用串行逻辑。

\textbf{正确性验证：}

通过恒等式$\text{FFT}(\text{FFT}(X)) = N \cdot X$检验。对所有测试配置（$N=2^{16}, 2^{18}, 2^{20}$，$P=1,2,4,8,16$），计算误差均在$10^{-16}$量级，与串行版本完全一致，验证了并行实现的正确性。

\subsection{数据打包优化}

\textbf{实现方法：}

在数据交换阶段应用\texttt{MPI\_Pack}/\texttt{MPI\_Unpack}进行通信优化。打包版本先将待发送的$2N/P$个双精度数通过\texttt{MPI\_Pack}序列化到连续字节缓冲区，再以\texttt{MPI\_PACKED}数据类型进行\texttt{MPI\_Sendrecv}传输，接收方通过\texttt{MPI\_Unpack}反序列化。非打包版本直接以\texttt{MPI\_DOUBLE}类型传输原始数组。两版本使用统一的\texttt{use\_pack}标志切换，除数据交换步骤外其余逻辑完全一致。

需要指出，本实现中交换的数据本身已是连续内存块，因此\texttt{MPI\_Pack}理论上并不必要（打包/解包引入了额外的内存拷贝开销）。实际观察到的性能差异反映了MPI实现层面对\texttt{MPI\_PACKED}与\texttt{MPI\_DOUBLE}数据类型的不同内部处理路径。

\subsection{性能分析}

\subsubsection{不同问题规模(N)和并行规模(P)下的运行时间}

表\ref{tab:runtime_results}展示了非打包版本在不同N和P组合下的原始运行时间。注意：不同N使用了不同的迭代次数（N=$2^{16}$: 100次，N=$2^{18}$: 20次，N=$2^{20}$: 5次），因此绝对时间仅在同一N行内可比。加速比分析使用P=1的MPI版本作为基准（非纯串行版本）。

\begin{table}[h!]
\centering
\caption{不同问题规模(N)和并行规模(P)下的运行时间(单位：秒)}
\label{tab:runtime_results}
\begin{tabular}{@{}lccccc@{}}
\toprule
\multirow{2}{*}{问题规模(N)} & \multicolumn{5}{c}{并行规模(进程数P)} \\
\cmidrule(l){2-6}
                              & P=1(基准) & P=2 & P=4 & P=8 & P=16 \\ \midrule
$N = 2^{16}$         & 2.237 & 5.282 & 6.357 & 2.920 & 1.571 \\
$N = 2^{18}$        & 2.023 & 6.144 & 6.259 & 2.801 & 1.386 \\
$N = 2^{20}$       & 2.362 & 8.530 & 8.866 & 4.103 & 1.686 \\
\bottomrule
\end{tabular}
\end{table}

\subsubsection{加速比与并行效率分析}

表\ref{tab:speedup}展示了加速比$S_p = T_1 / T_p$和并行效率$E_p = S_p / P$。

\begin{table}[h!]
\centering
\caption{不同配置下的加速比($S_p$)与并行效率($E_p$)}
\label{tab:speedup}
\begin{tabular}{@{}lccccc@{}}
\toprule
\multirow{2}{*}{N} & 指标 & \multicolumn{4}{c}{并行规模(进程数P)} \\
\cmidrule(l){3-6}
                    &      & P=2 & P=4 & P=8 & P=16 \\ \midrule
\multirow{2}{*}{$2^{16}$}
                    & $S_p$ & 0.42 & 0.35 & 0.77 & 1.42 \\
                    & $E_p$ & 0.21 & 0.09 & 0.10 & 0.09 \\ \addlinespace
\multirow{2}{*}{$2^{18}$}
                    & $S_p$ & 0.33 & 0.32 & 0.72 & 1.46 \\
                    & $E_p$ & 0.16 & 0.08 & 0.09 & 0.09 \\ \addlinespace
\multirow{2}{*}{$2^{20}$}
                    & $S_p$ & 0.28 & 0.27 & 0.58 & 1.40 \\
                    & $E_p$ & 0.14 & 0.07 & 0.07 & 0.09 \\
\bottomrule
\end{tabular}
\end{table}

\textbf{加速比分析：}

\begin{itemize}
    \item P=2和P=4时加速比均小于1（严重减速），总执行时间远超P=1。小规模并行下，进程间通信开销（MPI\_Sendrecv数据交换 + MPI\_Alltoallv输出重分布）远超计算并行化带来的收益。
    \item P=8时加速比为0.58--0.77，虽仍低于1但明显改善，说明通信开销占比开始下降。
    \item P=16时加速比达到1.40--1.46，首次超过基准，也是本次实验中唯一实现正向加速的配置。但考虑到投入了16倍的计算资源，1.4--1.5倍的加速比远低于线性加速的理想预期。
    \item 加速比随N变化不大（P=16时1.40--1.46），因为本实现中每进程的计算量与通信量均随$N/P$等比增长，计算通信比近似恒定。
\end{itemize}

\textbf{并行效率分析：}

并行效率$E_p = S_p / P$在所有配置下均处于0.07--0.21的低水平，说明投入的并行资源绝大部分被通信开销消耗。$E_p$呈近似$1/P$的衰减态势（P=2时0.14--0.21，P=16时0.09），表明当前实现的可扩展性存在结构性瓶颈。

\textbf{性能局限性的深层分析（与理论预期的差距）：}

本实现中通信复杂度为$O((N/P) \cdot \log_2 N)$且全程需要$2 \log_2 N$次全局通信（以$N=2^{20}$为例，单次FFT需40次MPI集体操作）。相比经典高性能并行FFT实现（如转置法、六步FFT等）通常仅在少数阶段进行全局数据重分布，本实验实现由于在每个FFT阶段都执行全局通信，因此通信次数明显偏多，导致整体通信开销较高。造成这一差距的根本原因有两点：

\begin{enumerate}
    \item \textbf{每阶段双通信模式：}本实现在每个FFT阶段都执行MPI\_Sendrecv（数据交换）和MPI\_Alltoallv（输出重分布），无法利用「前$\log_2(N/P)$个阶段蝶形对完全位于进程内部」这一性质。这是因为串行代码采用按频率抽取（DIF）变体，所有蝶形对均跨越数组前半部分与后半部分的边界，而这两部分始终分布在不同的进程组中（进程$0\sim P/2-1$持有前半部分，$P/2\sim P-1$持有后半部分）。因此即使是最早的阶段（$mj=1$）也需要跨进程通信。
    \item \textbf{通信-计算无法重叠：}当前MPI\_Sendrecv和MPI\_Alltoallv均为阻塞操作，通信期间进程闲置等待，无法利用非阻塞通信实现计算与数据传输的重叠。
\end{enumerate}

\textbf{优化路径分析（不改变报告既有数据，仅供讨论）：}

若对算法进行以下结构性改进，理论上可将P=16时的加速比从1.4倍提升至接近线性加速：

\begin{itemize}
    \item \textbf{采用转置法（Transpose FFT）：}将$N$点FFT视为二维$\sqrt{N} \times \sqrt{N}$（或更一般地$N/P \times P$）网格，先在各行做局部FFT，通过一次MPI\_Alltoall转置，再在各列做局部FFT。全程仅需1次Alltoall（前向FFT），相比本实现的每阶段双通信（$2 \log_2 N = 40$次），通信次数从$O(\log N)$降至$O(1)$。
    \item \textbf{引入按时间抽取（DIT）变体：}DIT变体的蝶形配对距离为$mj$（而非固定$N/2$），前$\log_2(N/P)$个阶段（$mj \leq N/P$）的所有蝶形对完全位于进程内部，无需任何通信。仅最后$\log_2 P$个阶段需要跨进程通信（二进制交换法，每阶段1次Sendrecv），总通信阶段数从$\log_2 N$降至$\log_2 P$（当$P=16$时从20降至4）。
    \item \textbf{使用非阻塞通信：}以MPI\_Isend/MPI\_Irecv替代MPI\_Sendrecv，在等待数据传输完成的同时执行不依赖通信数据的局部计算，实现计算与通信的重叠。
    \item \textbf{消除重复计算：}当前实现中进程$r$与其伙伴$r \oplus P/2$计算了完全相同的蝶形对（两者交换后拥有相同的数据），导致计算量翻倍。重构计算分配可使总计算量减半。
\end{itemize}

本实验的核心价值在于：通过真实数据定量地展示了通信设计对并行程序性能的决定性影响——即使算法在数值上正确，不当的通信模式也可能使并行化失去实际意义。这一反面案例恰好印证了并行计算课程的核心原则：\textbf{并行效率取决于通信最小化，而非仅并行化本身}。

\subsubsection{数据打包对性能的影响}

表\ref{tab:pack_comparison}展示了所有并行规模下打包与非打包版本的完整性能对比。

\begin{table}[h!]
\centering
\caption{数据打包对运行时间的影响(单位：秒)}
\label{tab:pack_comparison}
\begin{tabular}{@{}lcccccccc@{}}
\toprule
\multirow{2}{*}{N} & \multicolumn{2}{c}{P=2} & \multicolumn{2}{c}{P=4} & \multicolumn{2}{c}{P=8} & \multicolumn{2}{c}{P=16} \\
\cmidrule(l){2-3} \cmidrule(l){4-5} \cmidrule(l){6-7} \cmidrule(l){8-9}
                   & 非打包 & 打包   & 非打包 & 打包   & 非打包 & 打包   & 非打包 & 打包 \\ \midrule
$2^{16}$           & 5.282 & 5.609 & 6.357 & 6.456 & 2.920 & 2.272 & 1.571 & 1.553 \\
$2^{18}$           & 6.144 & 6.526 & 6.259 & 6.713 & 2.801 & 2.285 & 1.386 & 1.384 \\
$2^{20}$           & 8.530 & 8.927 & 8.866 & 8.873 & 4.103 & 4.213 & 1.686 & 1.683 \\
\bottomrule
\end{tabular}
\end{table}

从表\ref{tab:pack_comparison}可以得出以下结论：

\begin{itemize}
    \item \textbf{P=2,4（打包负面影响）：}在所有N下打包版本均比非打包版本慢。P=2时打包带来约5--6\%的性能下降，P=4时下降1--7\%。这是因为在进程数较少时每条消息较大（$2N/P$个双精度数），数据本身已是连续内存块，\texttt{MPI\_Pack/Unpack}引入的额外内存拷贝开销是纯粹的负担，没有任何通信效率补偿。
    \item \textbf{P=8（打包条件性获益）：}在N=$2^{16}$和$2^{18}$下打包版本分别快22\%和18\%，但在N=$2^{20}$时反而慢2.7\%。这一非单调行为说明\texttt{MPI\_PACKED}与\texttt{MPI\_DOUBLE}在MPI内部有不同的处理路径，在特定消息大小范围内（本实验中约32--128 KB，对应$N/P=2^{13}\sim 2^{15}$个复数）打包路径可能有更好的内部缓冲策略。当消息继续增大（N=$2^{20}$, P=8时约256 KB），打包/解包的计算开销重新超过内部优化收益。
    \item \textbf{P=16（打包无明显影响）：}所有N下打包与非打包性能差异在2\%以内，处于测量误差范围。此时每条消息很小（$N/P=2^{12}\sim 2^{16}$个复数，即8--128 KB），打包/解包的CPU开销与MPI内部处理差异的收益大致抵消。
    \item \textbf{总体结论：}对于本实验中的连续数据交换场景，MPI\_Pack/Unpack并非普适优化手段。其效果高度依赖于进程数（决定消息大小）和MPI实现细节，在部分配置下可能带来提升，也可能引入额外开销。对于非连续数据或需自定义数据布局的场景，打包技术可能有更大的实际价值。
\end{itemize}

\subsection{内存消耗分析(Valgrind Massif)}

使用Valgrind Massif工具（带\texttt{--stacks=yes}参数）对N=4096 ($2^{12}$)的并行FFT程序进行内存分析。

\begin{figure}[h!]
\centering
\caption{P=1时内存消耗图 (ms\_print输出, N=4096)}
\label{fig:massif_p1}
\begin{verbatim}
    MB
3.610^   #
     |   #@:@@:::::::::::::::::@:::::@:::@::::::::::@@::::@::
     |   #@:@@:::::::::::::: ::@: :::@:::@::::::::::@ ::::@::
     |  @#@:@@:::::::::::::: ::@: :::@:::@::::::::::@ ::::@::
     |  @#@:@@:::::::::::::: ::@: :::@:::@::::::::::@ ::::@::
     |  @#@:@@:::::::::::::: ::@: :::@:::@::::::::::@ ::::@::
     |  @#@:@@:::::::::::::: ::@: :::@:::@::::::::::@ ::::@::
     |  @#@:@@:::::::::::::: ::@: :::@:::@::::::::::@ ::::@::
     |  @#@:@@:::::::::::::: ::@: :::@:::@::::::::::@ ::::@::
     |  @#@:@@:::::::::::::: ::@: :::@:::@::::::::::@ ::::@::
     |  @#@:@@:::::::::::::: ::@: :::@:::@::::::::::@ ::::@::
     |  @#@:@@:::::::::::::: ::@: :::@:::@::::::::::@ ::::@::
     |  @#@:@@:::::::::::::: ::@: :::@:::@::::::::::@ ::::@::
     | :@#@:@@:::::::::::::: ::@: :::@:::@::::::::::@ ::::@::
     | :@#@:@@:::::::::::::: ::@: :::@:::@::::::::::@ ::::@::
   0 +------------------------------------------------------>Mi
     0                                                  708.6
Number of snapshots: 59
Detailed snapshots: [2, 3 (peak), 5, 7, 8, 25, 30, 34, 47, 54]
\end{verbatim}
\end{figure}

选择较小的N是为了使Valgrind（会使程序运行速度减慢10--50倍）能够在合理时间内完成采集。图\ref{fig:massif_p1}和图\ref{fig:massif_p4}分别展示了P=1和P=4时\texttt{ms\_print}生成的内存消耗图。图中x轴为程序运行时间（单位：百万指令Mi），y轴为内存消耗量（单位：MB），\texttt{\#}标记峰值快照，\texttt{@}标记详细快照，\texttt{:}标记普通快照。

\begin{figure}[h!]
\centering
\caption{P=4时内存消耗图 (ms\_print输出, N=4096, 进程47815)}
\label{fig:massif_p4}
\begin{verbatim}
    MB
3.532^                                                :
     |    #::::::::::::::::::::::::::::::::::::::::::::@:::@:
     |    #: ::: :::::::::: :::::::::: :: : ::::::: :::@: :@:
     |   @#: ::: :::::::::: :::::::::: :: : ::::::: :::@: :@:
     |   @#: ::: :::::::::: :::::::::: :: : ::::::: :::@: :@:
     |   @#: ::: :::::::::: :::::::::: :: : ::::::: :::@: :@:
     |   @#: ::: :::::::::: :::::::::: :: : ::::::: :::@: :@:
     |   @#: ::: :::::::::: :::::::::: :: : ::::::: :::@: :@:
     |   @#: ::: :::::::::: :::::::::: :: : ::::::: :::@: :@:
     |   @#: ::: :::::::::: :::::::::: :: : ::::::: :::@: :@:
     |   @#: ::: :::::::::: :::::::::: :: : ::::::: :::@: :@:
     |   @#: ::: :::::::::: :::::::::: :: : ::::::: :::@: :@:
     |  :@#: ::: :::::::::: :::::::::: :: : ::::::: :::@: :@:
     | ::@#: ::: :::::::::: :::::::::: :: : ::::::: :::@: :@:
     | ::@#: ::: :::::::::: :::::::::: :: : ::::::: :::@: :@:
   0 +------------------------------------------------------>Mi
     0                                                  524.3
Number of snapshots: 53
Detailed snapshots: [3, 4 (peak), 46, 49]
\end{verbatim}
\end{figure}

表\ref{tab:memory}总结了峰值快照的关键内存数据。

\begin{table}[h!]
\centering
\caption{峰值快照内存消耗数据 (N=4096, \texttt{--stacks=yes})}
\label{tab:memory}
\begin{tabular}{@{}lccccc@{}}
\toprule
配置 & 快照编号 & 时刻(Mi) & 总内存(B) & 有用堆(B) & 栈内存(B) \\ \midrule
P=1 & 3 (peak) & 42.66 & 3,785,512 & 3,603,619 & 1,792 \\
P=4 (进程1) & 4 (peak) & 46.66 & 3,701,712 & 3,516,891 & 1,792 \\
\bottomrule
\end{tabular}
\end{table}

\textbf{内存分配来源分析：}

峰值时刻的Massif分配树显示：
\begin{itemize}
    \item P=1：有用堆内存中42.69\%（约1.54 MB）来自MPI内部通信缓冲区的内存池函数\texttt{opal\_free\_list\_grow\_st}（属于OpenMPI的btl/base组件），10.18\%来自OpenMPI其他初始化分配（变量注册、进程管理等），仅约6.41\%来自程序自身的数据结构分配。
    \item P=4：43.83\%来自同一MPI通信缓冲池，14.93\%来自其他MPI分配，与P=1的分配模式高度相似。
    \item 两个配置的栈内存消耗均为1,792 B（约1.75 KB），证实了程序几乎完全使用堆内存，调用栈深度浅、栈帧尺寸小，符合FFT迭代算法的预期。
\end{itemize}

\textbf{内存消耗变化趋势分析：}

从图\ref{fig:massif_p1}和图\ref{fig:massif_p4}的内存消耗曲线可以看出：

\begin{itemize}
    \item \textbf{初始化阶段（图表左侧上升段）：}程序在MPI初始化和首轮内存分配阶段快速消耗内存。P=1在约42.66 Mi指令处即达到峰值，P=4稍晚（约46.66 Mi），这源于多进程协调（握手、barrier同步等）引入的额外指令开销。
    \item \textbf{稳态阶段（图表中间平坦段）：}内存消耗在FFT计算过程中保持几乎完全恒定。大量\texttt{:}字符构成的平坦区域说明程序在初始化后不再有显著的净堆分配——数据数组在初始化时一次分配并在多次FFT迭代中反复复用，符合高效内存使用的设计原则。偶尔出现的\texttt{@}标记（详细快照）可能捕获了\texttt{step\_mpi}内部临时向量的分配与释放，但净内存变化极小。
    \item \textbf{进程间对比：}P=1的单进程峰值总内存为3.61 MB（3,785,512 B）。P=4的单个进程峰值内存为3.53 MB（3,701,712 B），略低于P=1，这是因为数据块大小从N=4096（P=1）降至N/P=1024（P=4），每进程的数据数组尺寸减小。然而4个进程的内存消耗总和约为$4 \times 3.53 = 14.12$ MB，约等于P=1的4倍，体现了分布式内存模型的典型特征：总内存消耗随进程数线性增长，峰值内存主要受MPI运行时框架（约占50\%以上）而非应用数据的控制。
    \item \textbf{内存组成：}MPI运行时框架（OpenMPI）占据了峰值内存的50\%以上，包括通信缓冲区（\texttt{opal\_free\_list\_grow\_st}）和各种初始化分配。程序自身的FFT数据数组及临时向量仅占总内存的小部分。这意味着对于较小的问题规模（如N=4096），减少N并不能显著降低每进程的内存消耗，因为MPI框架的固定开销构成了"内存地板"。
\end{itemize}

\section{实验总结}

本实验以MPI对快速傅里叶变换（DIF变体）进行了并行化实现，分别测试了基础通信版本与MPI\_Pack/Unpack打包优化版本的性能，并利用Valgrind Massif分析了内存消耗特征。主要结论如下：

\begin{enumerate}
    \item \textbf{并行加速效果：}在P=16时可获得约1.4倍加速比，但P=2,4,8时均出现不同程度的减速（加速比0.28--0.77）。P=16时1.4倍的加速比远低于投入16倍计算资源的线性加速理想值，根本原因在于当前实现的通信模式：由于当前实现采用固定块状分布，并在每阶段维持“前半区数据位于低编号进程、后半区数据位于高编号进程”的布局方式，因此每个DIF阶段都需要跨进程交换蝶形输入数据。这属于当前数据布局与通信策略带来的结果，而非DIF算法本身的必然特性。每个FFT阶段（共$\log_2 N$个）都需要两次全局集体通信（Sendrecv + Alltoallv），导致通信开销完全吞噬了计算并行化收益。
    \item \textbf{通信模式的教训：}标准并行FFT可通过转置法或DIT变体将通信阶段数从$\log_2 N$降至$\log_2 P$，本实验的实测数据恰好从反面验证了这一原则的必要性——不当的通信设计足以使并行化失去实际价值。
    \item \textbf{数据打包评估：}MPI\_Pack/Unpack的效果高度依赖于进程数和消息大小，并非普适优化：P=2,4时一致劣化5--7\%，P=8时条件性获益18--22\%，P=16时无明显影响。对于连续内存块的传输场景，打包操作的额外内存拷贝是纯粹的负担。
    \item \textbf{内存消耗特征：}Valgrind Massif分析表明每进程峰值内存约3.5--3.6 MB（N=4096），其中MPI运行时框架占50\%以上。总内存消耗随进程数线性增长，程序在初始化后维持恒定内存占用。
\end{enumerate}

\end{document}
